\documentclass[11pt, letterpaper]{article}
\usepackage[margin=1in]{geometry}
\usepackage{titlesec}
\usepackage{enumitem}

% Formatting titles
\titleformat{\section}{\large\bfseries}{\thesection}{1em}{}
\titlespacing*{\section}{0pt}{1.5ex plus .5ex minus .2ex}{0.5ex plus .2ex}

\title{\vspace{-2cm}\textbf{EC519 Project Proposal}}
\vspace{-1cm}
\author{Mimi Paule}
\date{}

\begin{document}

\maketitle
\thispagestyle{empty} % Remove page number on the first page
\vspace{-1cm}

\section*{Project Title: Cross-Linguistic Phoneme Analysis via Short-Time LPC Formant Tracking and Dynamic Time Warping}
\vspace{0.3cm}
\section{Introduction and Objectives}
This project investigates the acoustic differences in vocal tract dynamics across different languages by analyzing words that sound similar phonetically. 
The primary objective is to apply Short-Time Linear Predictive Coding (LPC) to track changing formant frequencies ($F_1, F_2, F_3$) over time. 
By using these formant trajectories as acoustic features, we will then use Dynamic Time Warping (DTW) to calculate the "acoustic distance" between the pronunciations. 
This project will utilize core course concepts of Short-Time Analysis, Machine Deconvolution, and Speech Recognition algorithms into a single application, quantifying how identical words physically differ across native linguistic backgrounds.

\section{Proposed Algorithms}
\begin{enumerate}[noitemsep,topsep=0pt]
    \item \textbf{Dynamic Feature Extraction (Short-Time LPC):} The audio signals will be partitioned into overlapping frames using a Hamming window. 
    For each frame, the vocal tract will be modeled as an all-pole IIR filter. 
    While standard built-in libraries in Python will be utilized to calculate the raw LPC coefficients computationally efficiently, further root-finding algorithms will be manually implemented to convert those polynomial coefficients into exact numerical formant frequencies.
    \item \textbf{Sequence Alignment (DTW):} The resulting formant tracks will be fed into a Dynamic Time Warping algorithm. DTW will align the dynamic sequence lengths of the pronounced words to compute an optimal match and return a scalar distance metric that will indicate which languages share closer acoustic articulation geometries.
\end{enumerate}

\section{Data Description}
The project will utilize actual speech data consisting of a list of self-recorded words spoken by myself in my native languages (English and Japanese) as well as a language I am currently learning (Arabic). 
This dataset will be collected using a standardized microphone setup across all three languages to ensure consistent recording conditions. 
%The data will undergo systematic preprocessing including downsampling to a standardized rate and framing/windowing prior to the short-time LPC analysis.

\section{Expected Outcomes}
The expected outcome is an application capable of tracking dynamic formants during active speech and quantifying the linguistic differences through sequence alignment. 
The results will be visualized using temporal formant tracks overlaid onto spectrograms, as well as a computed distance matrix generated by the DTW algorithm. 
%This visual and statistical analysis will validate how well short-time LPC features combined with DTW can capture subtle, cross-linguistic acoustic differences in equivalent words.

\section{Timeline}
\begin{itemize}[noitemsep,topsep=0pt]
    \item \textbf{Data Collection \& Preprocessing:} Record the equivalent word datasets across English, Japanese, and Arabic; establish the windowing/STFT analysis pipeline.
    \item \textbf{Dynamic LPC Implementation:} Implement the Short-Time LPC algorithm across sliding frames to estimate $F_1$ and $F_2$ formants, extracting their temporal trajectories.
    \item \textbf{DTW Alignment \& Analysis:} Implement the Dynamic Time Warping logic to calculate the acoustic distances between the extracted formant tracks, generate spectrogram overlays, and write up the final report.
\end{itemize}

\end{document}
